\section{How to Evaluate Honestly}
\label{sec:pitfalls}

The hardest part of quantum-enabled learning is not building the model but trusting
the comparison. The wider machine-learning community has learned this lesson the hard
way, in reinforcement learning where seeds and implementation details dominate
reported gains~\cite{henderson2018matters}, in generative modeling where a fair
large-scale comparison erased claimed advantages~\cite{lucic2018gans}, and in
benchmarking practice more broadly~\cite{pineau2021reproducibility,dehghani2021lottery}.
This section reproduces two pitfalls that manufacture false advantages, shows how to
remove them, and distills a checklist. These are not hypothetical: each one produced a
misleading signal in our own early experiments before we corrected it.

\subsection{Pitfall 1: scale-coupled normalization}
A node-potential target has units of hops, and those units grow with the graph. If we
normalize the target by a global graph constant such as the diameter,
$y_i=d_G(i,v_d)/\mathrm{diam}(G)$, and then pool errors across graphs of different
size, large graphs dominate the average and the apparent global-versus-local gap
appears to explode with diameter. Fig.~\ref{fig:pitfall} shows the effect: under raw
hop units the gap grows with a slope of $\cRawSlope$ per unit diameter and reaches
$\cRawMax$ on the largest graph, whereas under a scale-free per-source normalization
(by eccentricity) the same gap grows with a slope of only $\cSFSlope$, roughly
$\cSlopeRatio\times$ flatter, and stays near $\cSFMax$.

The honest reading separates two things. A genuine, modest gap remains under the
scale-free metric, because a fixed $K$-hop GCN really does fall behind as the diameter
grows. What is spurious is the \emph{magnitude and steepness} of the trend under
raw units, which is an artifact of letting the metric scale with the problem. The fix
is to normalize per source to a scale-free quantity before pooling heterogeneous
graphs, and to report a scale-free summary such as the AURC, in the spirit of
standardized robustness protocols~\cite{hendrycks2019corruptions}.

\begin{figure}[t]
\centering
\resizebox{\columnwidth}{!}{%
\begin{tikzpicture}
\begin{axis}[
  width=8.4cm, height=5.7cm,
  xlabel={graph diameter}, ylabel={apparent global-vs-local gap},
  xmin=5, xmax=19, ymin=-0.06, ymax=0.7,
  legend pos=north west, legend cell align=left,
  grid=major, grid style={gray!18},
  tick label style={font=\footnotesize}, label style={font=\small},
  legend style={font=\footnotesize, fill=white, fill opacity=0.85, draw=gray!40},
  axis line style={gray!50},
]
% scale-free (eccentricity) normalization: flat
\addplot[steel, very thick, mark=*, mark size=1.4pt, name path=sf] coordinates {
  (6,-0.020)(8,-0.023)(10,0.018)(12,0.031)(14,0.044)(16,0.050)(18,0.053)};
\addlegendentry{scale-free per source (fix)}
% raw normalization: spurious runaway
\addplot[amber!80!black, very thick, mark=square*, mark size=1.4pt, name path=raw]
  coordinates {
  (6,-0.024)(8,0.019)(10,0.161)(12,0.278)(14,0.483)(16,0.589)(18,0.632)};
\addlegendentry{raw units (pitfall)}
% shade the spurious gap between the two
\addplot[amber!30, opacity=0.45] fill between[of=raw and sf];
\node[font=\footnotesize\bfseries, text=amber!70!black, align=center]
  at (axis cs:14.2,0.30) {scale\\artifact};
\end{axis}
\end{tikzpicture}}
\caption{Exp~C1, the normalization pitfall. Pooling errors in raw hop units across
graphs of growing diameter manufactures a gap that explodes to $\cRawMax$ (slope
$\cRawSlope$). A scale-free per-source normalization leaves a small, genuine gap
(slope $\cSFSlope$, about $\cSlopeRatio\times$ flatter): the real long-range deficiency
of a local operator, without the scale artifact (shaded).}
\label{fig:pitfall}
\end{figure}


\subsection{Pitfall 2: single-seed reporting}
Quantum-inspired models in our study are visibly noisier than their classical
counterparts. On the $66$-satellite shell over eight seeds, the quantum-inspired walk
ranges widely while the heat kernel is tight: $\cTwoQw\pm\cTwoQwStd$ versus
$\cTwoHeat\pm\cTwoHeatStd$ (the local GCN sits at $\cTwoGcn\pm\cTwoGcnStd$). A single
seed could therefore show the quantum-inspired operator beating, matching, or losing
to the classical one, as Fig.~\ref{fig:seeds} makes visible. Reporting one run is not a
small sin here; it can invert the conclusion. The remedy is to report mean and standard
deviation over several seeds and to treat any ranking that lies inside the combined
spread as a tie.

\begin{figure}[t]
\centering
\resizebox{\columnwidth}{!}{%
\begin{tikzpicture}
\begin{axis}[
  width=8.4cm, height=5.7cm,
  xlabel={seed}, ylabel={MAE (lower is better)},
  xmin=-0.5, xmax=7.5, ymin=0.04, ymax=0.13,
  xtick={0,1,2,3,4,5,6,7},
  legend pos=north east, legend cell align=left,
  grid=major, grid style={gray!18},
  tick label style={font=\footnotesize}, label style={font=\small},
  legend style={font=\scriptsize, fill=white, fill opacity=0.85, draw=gray!40},
  axis line style={gray!50},
  clip=false,
]
% mean +/- std bands (drawn behind the markers)
\fill[amber!22] (axis cs:-0.5,0.073) rectangle (axis cs:7.5,0.097);
\fill[steel!22] (axis cs:-0.5,0.059) rectangle (axis cs:7.5,0.069);
\draw[amber!80!black, dashed, thick] (axis cs:-0.5,0.085) -- (axis cs:7.5,0.085);
\draw[steel, dashed, thick] (axis cs:-0.5,0.064) -- (axis cs:7.5,0.064);
% per-seed markers
\addplot[only marks, mark=square*, amber!80!black, mark size=2pt] coordinates {
  (0,0.0959)(1,0.0705)(2,0.1046)(3,0.0736)(4,0.0691)(5,0.0895)(6,0.0921)(7,0.0823)};
\addlegendentry{QI (wide spread)}
\addplot[only marks, mark=*, steel, mark size=2pt] coordinates {
  (0,0.0658)(1,0.0586)(2,0.0642)(3,0.0638)(4,0.0542)(5,0.0687)(6,0.0616)(7,0.0726)};
\addlegendentry{Heat (tight)}
\end{axis}
\end{tikzpicture}}
\caption{Exp~C2, the single-seed pitfall ($66$ satellites, $8$ seeds). Dashed lines mark
the per-operator means and shaded bands the $\pm$ one standard-deviation ranges. The
quantum-inspired walk (QI) scatters widely while the heat kernel is tight; a single seed
could show QI beating, matching, or losing to heat. Only the distribution,
$\cTwoQw\pm\cTwoQwStd$ versus $\cTwoHeat\pm\cTwoHeatStd$, is trustworthy.}
\label{fig:seeds}
\end{figure}


\subsection{Pitfall 3: survivorship in the denominator}
A subtler trap appears whenever some instances can fail to produce a usable answer, for
example a routing decoder that does not deliver every flow, or a method scored only where
it converged. Averaging the error over the \emph{surviving} instances rewards a method
that quietly drops its hardest cases: the harder the case it abandons, the better its
conditional mean looks. The fix is to fix the denominator, every instance counts and a
dropped case is scored at its true (failed) cost rather than omitted. The same logic
applies to seeds and configurations: report all runs, not only the ones that trained
cleanly.

\subsection{Pitfall 4: the capacity confound in quantum comparisons}
Quantum and classical models are easy to compare unfairly because the quantum knob often
moves capacity at the same time. Adding qubits typically widens the classical pre- or
post-processing as well, so a quantum advantage measured at matched qubit count can be a
plain capacity advantage in disguise. The remedy is to match the resource that actually
carries the capacity, the bottleneck width and total parameter count, not the qubit
count. We apply this directly: the multi-time quantum-walk variant in
Section~\ref{sec:casestudy} is given more parameters and is reported with its larger
count, so the reader can see it does not buy an improvement despite the extra capacity.

\subsection{A reusable checklist}
\begin{enumerate}
\item \textbf{Param-matched operators.} Vary only the propagation operator; hold depth,
width, and weight shapes fixed.
\item \textbf{Multiple seeds.} Report mean $\pm$ std; a ranking inside the std is not a
result.
\item \textbf{Scale-free metrics.} Normalize per source and report a scale-free summary
(AURC); never pool raw units across sizes.
\item \textbf{Honest denominators.} Count every instance; beware survivorship.
\item \textbf{Design sweep before a verdict.} Probe the obvious variants so that ``it
did not help'' is a statement about the mechanism, not one implementation.
\end{enumerate}
